% --- عنوان اصلی ---
  \node[text=white, font=\fontsize{32}{38}\selectfont\bfseries,
    anchor=center, align=center] at (10.5,17.5)
    {\rl{تحلیل جامع انقلاب فرانسه}};%
  % --- زیرعنوان ---
  \node[text=AccentGold, font=\fontsize{18}{24}\selectfont\bfseries,
    anchor=center, align=center] at (10.5,15.5)
    {\rl{۱۷۸۹–۱۷۹۹ و تحولات پس از آن}};%
  % --- شعار ---
  \node[text=white, font=\fontsize{14}{18}\selectfont\itshape,
    anchor=center, align=center] at (10.5,12.2)
    {Libert\'{e}, \'{E}galit\'{e}, Fraternit\'{e}};%
  \node[text=NeutralMid, font=\fontsize{12}{16}\selectfont,
    anchor=center, align=center] at (10.5,11.4)
    {\rl{آزادی، برابری، برادری}};%
  % --- نماد باستیل ساده‌شده ---
  \begin{scope}[shift={(10.5,7)}, scale=0.8]
    % برج چپ
    \fill[NeutralMid, opacity=0.3] (-2,0) rectangle (-1,3);%
    \fill[NeutralMid, opacity=0.3] (-2.2,3) rectangle (-0.8,3.5);%
    % برج راست
    \fill[NeutralMid, opacity=0.3] (1,0) rectangle (2,3);%
    \fill[NeutralMid, opacity=0.3] (0.8,3) rectangle (2.2,3.5);%
    % دیوار وسط
    \fill[NeutralMid, opacity=0.2] (-1,0) rectangle (1,2);%
    % دروازه
    \fill[PrimaryDeep] (-0.4,0) arc (180:0:0.4) -- (0.4,0) -- cycle;%
    % شکست‌ها
    \draw[AccentGold, thick] (-1.5,2.5) -- (-1.2,1.5) -- (-1.6,0.5);%
    \draw[AccentGold, thick] (1.5,2.5) -- (1.3,1.5) -- (1.7,0.5);%
    \draw[AccentGold, thick, opacity=0.7] (0,2) -- (0.3,1) -- (-0.2,0.3);%
  \end{scope}
  % --- نام نویسنده ---
  \node[text=NeutralMid, font=\large, anchor=center] at (10.5,3.5)
    {\rl{مهدی سالم}};%
  % --- تاریخ ---
  \node[text=NeutralMid, font=\normalsize, anchor=center] at (10.5,2.5)
    {\rl{۱۴۰۴}};%
\end{tikzpicture}
\end{titlepage}

% ─────────────────────────────────────────────────────────────────
%   فهرست مطالب
% ─────────────────────────────────────────────────────────────────
\newpage
\tableofcontents
\newpage

% ═══════════════════════════════════════════════════════════════════
%   خلاصه‌ی اجرایی
% ═══════════════════════════════════════════════════════════════════
\needspace{6\baselineskip}%
\section*{خلاصه‌ی اجرایی}
\addcontentsline{toc}{section}{خلاصه‌ی اجرایی}

\begin{tcolorbox}[mybox, title={\rl{انقلاب فرانسه: مادر انقلاب‌های مدرن}}]
انقلاب فرانسه \concept{مادر انقلاب‌های مدرن} است که الگوی تحول سیاسی رادیکال را
برای دو قرن آینده تعریف کرد. این انقلاب در \concept{چهار فاز} اصلی قابل تحلیل
است و میراث آن از اعلامیه‌ی حقوق بشر تا ناسیونالیسم مدرن گسترده است.
\end{tcolorbox}

\begin{tcolorbox}[defbox, title={\rl{چهار فاز انقلاب}}]
\begin{center}
\begin{adjustbox}{max width=\linewidth}% FIXED: جلوگیری از سرریز
\begin{tabular}{>{\bfseries\color{white}}r
                >{\small}r
                >{\small}r
                >{\small}r}
\toprule
\rowcolor{PrimaryDeep}
\color{white}فاز & \color{white}دوره & \color{white}ماهیت & \color{white}نتیجه \\
\midrule
\rowcolor{PhaseLiberal!10}
\color{PhaseLiberal}لیبرال–مشروطه & ۱۷۸۹–۱۷۹۱ & اصلاحات از بالا & سلطنت مشروطه \\
\rowcolor{PhaseRadical!10}
\color{PhaseRadical}رادیکال–جمهوری & ۱۷۹۲–۱۷۹۴ & انقلاب توده‌ای & ترور و جمهوری \\
\rowcolor{PhaseThermidor!10}
\color{PhaseThermidor}ترمیدور و واکنش & ۱۷۹۴–۱۷۹۹ & تعدیل و بازگشت & دیرکتوار \\
\rowcolor{PhaseNapoleon!10}
\color{PhaseNapoleon}ناپلئونی & ۱۷۹۹–۱۸۱۵ & استبداد روشنگرانه & امپراتوری \\
\bottomrule
\end{tabular}
\end{adjustbox}
\end{center}
\end{tcolorbox}

\subsection*{میراث اصلی}
\begin{itemize}[nosep, rightmargin=1em]
  \item اعلامیه‌ی حقوق بشر و شهروند
  \item مفهوم \concept{ملت} به جای \concept{رعیت}
  \item سکولاریسم و جدایی دین از دولت
  \item ناسیونالیسم مدرن
  \item الگوی انقلاب رادیکال
\end{itemize}

% ═══════════════════════════════════════════════════════════════════
%   نقشه‌ی مفهومی انقلاب
% ═══════════════════════════════════════════════════════════════════
\newpage
\needspace{6\baselineskip}%
\section*{نقشه‌ی مفهومی انقلاب فرانسه}
\addcontentsline{toc}{section}{نقشه‌ی مفهومی}

\begin{center}
\begin{adjustbox}{max width=\linewidth}% FIXED: جلوگیری از سرریز
\begin{tikzpicture}[
  every node/.style={align=center, font=\small},
  root/.style={circle, draw=PrimaryDeep, fill=PrimaryDeep!15,
    line width=1.5pt, minimum size=2.8cm, font=\bfseries\large},
  branch/.style={rectangle, rounded corners=6pt, draw=#1, fill=#1!12,
    line width=1pt, minimum width=2.2cm, minimum height=1cm, font=\bfseries\small},
  leaf/.style={rectangle, rounded corners=3pt, draw=#1!60, fill=#1!6,
    minimum width=1.8cm, minimum height=0.7cm, font=\scriptsize},
  arr/.style={-{Stealth[length=5pt]}, thick, #1},
]
  % --- ریشه ---
  \node[root] (rev) {\rl{انقلاب}\\[2pt]\rl{فرانسه}\\[1pt]\rl{۱۷۸۹}};

  % --- شاخه‌های اصلی ---
  \node[branch=PhaseRadical, above right=2.5cm and 3cm of rev] (causes)
    {\rl{علل}};
  \node[branch=PhaseLiberal, right=4cm of rev] (phases)
    {\rl{فازها}};
  \node[branch=AccentTeal, below right=2.5cm and 3cm of rev] (legacy)
    {\rl{میراث}};
  \node[branch=AccentPurple, above left=2.5cm and 3cm of rev] (actors)
    {\rl{بازیگران}};
  \node[branch=AccentAmber, below left=2.5cm and 3cm of rev] (docs)
    {\rl{اسناد کلیدی}};

  % --- خطوط اصلی ---
  \draw[arr=PhaseRadical] (rev) -- (causes);
  \draw[arr=PhaseLiberal] (rev) -- (phases);
  \draw[arr=AccentTeal] (rev) -- (legacy);
  \draw[arr=AccentPurple] (rev) -- (actors);
  \draw[arr=AccentAmber] (rev) -- (docs);

  % --- برگ‌های علل ---
  \node[leaf=PhaseRadical, above=0.6cm of causes] (c1) {\rl{بحران مالی}};
  \node[leaf=PhaseRadical, right=0.4cm of causes] (c2) {\rl{نابرابری طبقاتی}};
  \node[leaf=PhaseRadical, above right=0.4cm and -0.5cm of causes] (c3) {\rl{روشنگری}};
  \node[leaf=PhaseRadical, below=0.6cm of causes] (c4) {\rl{قحطی ۱۷۸۸}};
  \draw[PhaseRadical!60, thin] (causes) -- (c1);
  \draw[PhaseRadical!60, thin] (causes) -- (c2);
  \draw[PhaseRadical!60, thin] (causes) -- (c3);
  \draw[PhaseRadical!60, thin] (causes) -- (c4);

  % --- برگ‌های فازها ---
  \node[leaf=PhaseLiberal, above=0.6cm of phases] (p1) {\rl{لیبرال ۸۹–۹۱}};
  \node[leaf=PhaseRadical, right=0.4cm of phases] (p2) {\rl{رادیکال ۹۲–۹۴}};
  \node[leaf=PhaseThermidor, below=0.6cm of phases] (p3) {\rl{ترمیدور ۹۴–۹۹}};
  \node[leaf=PhaseNapoleon, below right=0.4cm and -0.5cm of phases] (p4)
    {\rl{ناپلئونی ۹۹–۱۵}};
  \draw[PhaseLiberal!60, thin] (phases) -- (p1);
  \draw[PhaseRadical!60, thin] (phases) -- (p2);
  \draw[PhaseThermidor!60, thin] (phases) -- (p3);
  \draw[PhaseNapoleon!60, thin] (phases) -- (p4);

  % --- برگ‌های میراث ---
  \node[leaf=AccentTeal, below=0.6cm of legacy] (l1) {\rl{حقوق بشر}};
  \node[leaf=AccentTeal, right=0.4cm of legacy] (l2) {\rl{سکولاریسم}};
  \node[leaf=AccentTeal, below right=0.4cm and -0.5cm of legacy] (l3)
    {\rl{ناسیونالیسم}};
  \draw[AccentTeal!60, thin] (legacy) -- (l1);
  \draw[AccentTeal!60, thin] (legacy) -- (l2);
  \draw[AccentTeal!60, thin] (legacy) -- (l3);

  % --- برگ‌های بازیگران ---
  \node[leaf=AccentPurple, above=0.6cm of actors] (a1) {\rl{روبسپیر}};
  \node[leaf=AccentPurple, left=0.4cm of actors] (a2) {\rl{دانتون}};
  \node[leaf=AccentPurple, above left=0.4cm and -0.5cm of actors] (a3)
    {\rl{ناپلئون}};
  \node[leaf=AccentPurple, below=0.6cm of actors] (a4) {\rl{لافایت}};
  \draw[AccentPurple!60, thin] (actors) -- (a1);
  \draw[AccentPurple!60, thin] (actors) -- (a2);
  \draw[AccentPurple!60, thin] (actors) -- (a3);
  \draw[AccentPurple!60, thin] (actors) -- (a4);

  % --- برگ‌های اسناد ---
  \node[leaf=AccentAmber, below=0.6cm of docs] (d1) {\rl{اعلامیه‌ی حقوق بشر}};
  \node[leaf=AccentAmber, left=0.4cm of docs] (d2) {\rl{قانون اساسی ۹۱}};
  \node[leaf=AccentAmber, below left=0.4cm and -0.5cm of docs] (d3)
    {\rl{کد ناپلئون}};
  \draw[AccentAmber!60, thin] (docs) -- (d1);
  \draw[AccentAmber!60, thin] (docs) -- (d2);
  \draw[AccentAmber!60, thin] (docs) -- (d3);

\end{tikzpicture}
\end{adjustbox}
\end{center}

% ═══════════════════════════════════════════════════════════════════
%   اینفوگرافی خط زمان
% ═══════════════════════════════════════════════════════════════════
\newpage
\needspace{6\baselineskip}%
\section*{اینفوگرافی: خط زمان حوادث کلیدی}
\addcontentsline{toc}{section}{اینفوگرافی خط زمان}

\begin{center}
\begin{adjustbox}{max width=\linewidth}% FIXED: جلوگیری از سرریز
\begin{tikzpicture}[
  every node/.style={font=\scriptsize, align=center},
  event/.style={rectangle, rounded corners=4pt, draw=#1, fill=#1!10,
    minimum width=2cm, minimum height=0.7cm, font=\scriptsize\bfseries},
]
  % --- خط زمان اصلی ---
  \draw[PrimaryDeep, line width=2pt, -{Stealth[length=6pt]}]
    (0,0) -- (16,0);

  % علامت‌های سال
  \node[event=PhaseLiberal, above=0.8cm] at (1,0) (e1)
    {\rl{سقوط باستیل}\\[1pt]\rl{۱۴ ژوئیه ۱۷۸۹}};
  \draw[PhaseLiberal, thick] (1,0.1) -- (e1);

  \node[event=PhaseLiberal, below=0.8cm] at (2.8,0) (e2)
    {\rl{اعلامیه‌ی حقوق بشر}\\[1pt]\rl{۲۶ اوت ۱۷۸۹}};
  \draw[PhaseLiberal, thick] (2.8,-0.1) -- (e2);

  \node[event=PhaseLiberal, above=0.8cm] at (4.5,0) (e3)
    {\rl{قانون اساسی ۱۷۹۱}};
  \draw[PhaseLiberal, thick] (4.5,0.1) -- (e3);

  \node[event=PhaseRadical, below=0.8cm] at (6.2,0) (e4)
    {\rl{سقوط سلطنت}\\[1pt]\rl{۱۰ اوت ۱۷۹۲}};
  \draw[PhaseRadical, thick] (6.2,-0.1) -- (e4);

  \node[event=PhaseRadical, above=0.8cm] at (7.8,0) (e5)
    {\rl{اعدام لویی شانزدهم}\\[1pt]\rl{ژانویه ۱۷۹۳}};
  \draw[PhaseRadical, thick] (7.8,0.1) -- (e5);

  \node[event=PhaseRadical, below=0.8cm] at (9.5,0) (e6)
    {\rl{دوران ترور}\\[1pt]\rl{۱۷۹۳–۹۴}};
  \draw[PhaseRadical, thick] (9.5,-0.1) -- (e6);

  \node[event=PhaseThermidor, above=0.8cm] at (11,0) (e7)
    {\rl{سقوط روبسپیر}\\[1pt]\rl{ژوئیه ۱۷۹۴}};
  \draw[PhaseThermidor, thick] (11,0.1) -- (e7);

  \node[event=PhaseThermidor, below=0.8cm] at (12.5,0) (e8)
    {\rl{دیرکتوار}\\[1pt]\rl{۱۷۹۵–۹۹}};
  \draw[PhaseThermidor, thick] (12.5,-0.1) -- (e8);

  \node[event=PhaseNapoleon, above=0.8cm] at (14.5,0) (e9)
    {\rl{کودتای ناپلئون}\\[1pt]\rl{نوامبر ۱۷۹۹}};
  \draw[PhaseNapoleon, thick] (14.5,0.1) -- (e9);

  % --- راهنمای رنگ ---
  \node[event=PhaseLiberal] at (4,-3) (lg1) {\rl{فاز لیبرال}};
  \node[event=PhaseRadical, right=0.3cm of lg1] (lg2) {\rl{فاز رادیکال}};
  \node[event=PhaseThermidor, right=0.3cm of lg2] (lg3) {\rl{ترمیدور}};
  \node[event=PhaseNapoleon, right=0.3cm of lg3] (lg4) {\rl{ناپلئونی}};

\end{tikzpicture}
\end{adjustbox}
\end{center}

% ═══════════════════════════════════════════════════════════════════
%   طرح کلی تحولات فرانسه
% ═══════════════════════════════════════════════════════════════════
\newpage
\needspace{6\baselineskip}%
\section*{طرح کلی تحولات فرانسه (۱۷۸۹–۲۰۲۴)}
\addcontentsline{toc}{section}{طرح کلی تحولات فرانسه}

\begin{center}
\begin{adjustbox}{max width=\linewidth}% FIXED
\begin{tikzpicture}[
  every node/.style={font=\small, align=center},
  block/.style={rectangle, rounded corners=5pt, draw=#1, fill=#1!10,
    minimum width=4cm, minimum height=0.9cm, line width=1pt},
  arr/.style={-{Stealth[length=5pt]}, thick, NeutralDark},
  node distance=0.8cm,
]
  \node[block=PhaseRadical] (r1) {\rl{۱۷۸۹: انقلاب}};
  \node[block=PhaseLiberal, below=of r1] (r2) {\rl{۱۷۸۹–۱۷۹۲: سلطنت مشروطه}};
  \node[block=PhaseRadical, below=of r2] (r3) {\rl{۱۷۹۲–۱۸۰۴: جمهوری اول}};
  \node[block=PhaseNapoleon, below=of r3] (r4) {\rl{۱۸۰۴–۱۸۱۵: امپراتوری اول}};
  \node[block=EraAncient, below=of r4] (r5) {\rl{۱۸۱۴–۱۸۳۰: بازگشت بوربون‌ها}};
  \node[block=EraEarlyMod, below=of r5] (r6) {\rl{۱۸۳۰–۱۸۴۸: سلطنت ژوئیه}};
  \node[block=AccentTeal, below=of r6] (r7) {\rl{۱۸۴۸–۱۸۵۲: جمهوری دوم}};
  \node[block=AccentPurple, below=of r7] (r8) {\rl{۱۸۵۲–۱۸۷۰: امپراتوری دوم}};
  \node[block=EraModern, below=of r8] (r9) {\rl{۱۸۷۰–۱۹۴۰: جمهوری سوم}};
  \node[block=AccentRed, below=of r9] (r10) {\rl{۱۹۴۰–۱۹۴۴: ویشی}};
  \node[block=EraContemp, below=of r10] (r11) {\rl{۱۹۴۶–۱۹۵۸: جمهوری چهارم}};
  \node[block=PrimaryDeep, below=of r11] (r12) {\rl{۱۹۵۸–اکنون: جمهوری پنجم}};

  \foreach \i/\j in {r1/r2, r2/r3, r3/r4, r4/r5, r5/r6,
    r6/r7, r7/r8, r8/r9, r9/r10, r10/r11, r11/r12}{%
    \draw[arr] (\i) -- (\j);%
  }%
\end{tikzpicture}
\end{adjustbox}
\end{center}

% ═══════════════════════════════════════════════════════════════════
%   بخش اول: زمینه‌ها و علل انقلاب
% ═══════════════════════════════════════════════════════════════════
\newpage
\needspace{6\baselineskip}%
\section{زمینه‌ها و علل انقلاب}
\label{sec:causes}

\subsection{بحران ساختاری نظام کهن}
\label{sec:ancien}

\subsubsection{ساختار طبقاتی سه‌گانه}

جامعه‌ی فرانسه پیش از انقلاب بر اساس نظامی سه‌طبقه‌ای
(\lat{Ancien Régime}) سامان یافته بود که تناقض بنیادینی میان قدرت اقتصادی
و قدرت سیاسی ایجاد کرده بود.

\begin{center}
\begin{adjustbox}{max width=\linewidth}%
\begin{tabular}{rrrr}
\toprule
\rowcolor{PrimaryDeep}
\color{white}\bfseries طبقه & \color{white}\bfseries درصد جمعیت
  & \color{white}\bfseries ویژگی & \color{white}\bfseries امتیازات \\
\midrule
\rowcolor{BgWarm}
\textbf{اول: روحانیت} & ۰٫۵٪ (~۱۳۰٬۰۰۰) & کلیسای کاتولیک
  & معاف از مالیات، مالک ۱۰٪ زمین \\
\rowcolor{BgCool}
\textbf{دوم: اشراف} & ۱٫۵٪ (~۴۰۰٬۰۰۰) & نجبا، درباریان
  & معاف از مالیات، مالک ۲۵٪ زمین \\
\rowcolor{BgWarm}
\textbf{سوم: عوام} & ۹۸٪ (~۲۷ میلیون) & دهقانان، بورژوا، کارگران
  & تحمل همه‌ی مالیات‌ها \\
\bottomrule
\end{tabular}
\end{adjustbox}
\end{center}

\begin{tcolorbox}[wavebox, title={\rl{تناقض بنیادین}}]
\begin{center}
\begin{tikzpicture}[every node/.style={font=\small, align=center},
  block/.style={rectangle, rounded corners=4pt, draw=AccentRed,
    fill=AccentRed!8, minimum width=3.5cm, minimum height=0.7cm},
  arr/.style={-{Stealth[length=5pt]}, thick, AccentRed},
]
  \node[block] (a) {\rl{قدرت اقتصادی بورژوازی $\uparrow$}};
  \node[block, below=0.5cm of a] (b) {\rl{قدرت سیاسی اشراف: ثابت}};
  \node[block, below=0.5cm of b] (c) {\rl{تناقض ساختاری}};
  \node[block, below=0.5cm of c] (d) {\rl{بحران مشروعیت}};
  \draw[arr] (a) -- (b);
  \draw[arr] (b) -- (c);
  \draw[arr] (c) -- (d);
\end{tikzpicture}
\end{center}
\end{tcolorbox}

% ─── ۱.۲ بحران مالی ──────────────────────────────────────────
\needspace{6\baselineskip}%
\subsection{بحران مالی}
\label{sec:financial}

\subsubsection{ورشکستگی دولت}

دولت فرانسه در آستانه‌ی انقلاب با بحران مالی عمیقی دست‌وپنجه نرم می‌کرد.
هزینه‌ی جنگ استقلال آمریکا، اسراف دربار ورسای، و ساختار مالیاتی ناکارآمد
که طبقات ممتاز را معاف می‌کرد، کشور را به ورشکستگی نزدیک کرده بود.

\begin{center}
\begin{adjustbox}{max width=\linewidth}%
\begin{tabular}{rr}
\toprule
\rowcolor{PrimaryDeep}
\color{white}\bfseries منبع بحران & \color{white}\bfseries میزان \\
\midrule
\rowcolor{BgWarm} بدهی دولت (۱۷۸۸) & ۴ میلیارد لیور \\
\rowcolor{BgCool} هزینه‌ی سالانه‌ی بهره & ۵۰٪ بودجه \\
\rowcolor{BgWarm} کسری بودجه & ۱۲۶ میلیون لیور \\
\rowcolor{BgCool} هزینه‌ی جنگ استقلال آمریکا & ۱٫۳ میلیارد لیور \\
\bottomrule
\end{tabular}
\end{adjustbox}
\end{center}

\subsubsection{تلاش‌های ناموفق اصلاحی}

\begin{center}
\begin{adjustbox}{max width=\linewidth}%
\begin{tabular}{rrrr}
\toprule
\rowcolor{PrimaryDeep}
\color{white}\bfseries وزیر دارایی & \color{white}\bfseries سال
  & \color{white}\bfseries پیشنهاد & \color{white}\bfseries نتیجه \\
\midrule
\rowcolor{BgWarm} تورگو & ۱۷۷۴–۷۶ & اصلاحات لیبرال & برکناری \\
\rowcolor{BgCool} نکر & ۱۷۷۶–۸۱ & وام‌گیری & موقت \\
\rowcolor{BgWarm} کالون & ۱۷۸۳–۸۷ & مالیات بر اشراف & مخالفت اشراف \\
\rowcolor{BgCool} بریین & ۱۷۸۷–۸۸ & اصلاحات & شورش اشراف \\
\rowcolor{BgWarm} نکر (دوم) & ۱۷۸۸–۸۹ & مجلس طبقات & انقلاب \\
\bottomrule
\end{tabular}
\end{adjustbox}
\end{center}

% ─── ۱.۳ بحران ایدئولوژیک ────────────────────────────────────
\needspace{6\baselineskip}%
\subsection{بحران ایدئولوژیک: روشنگری}
\label{sec:enlightenment}

\subsubsection{متفکران کلیدی}

جنبش روشنگری (\lat{Les Lumières}) بنیادهای فکری انقلاب را فراهم کرد.
متفکرانی چون ولتر، مونتسکیو و روسو مشروعیت نظام کهن را به چالش کشیدند.

\begin{center}
\begin{adjustbox}{max width=\linewidth}%
\begin{tabular}{rrr}
\toprule
\rowcolor{PrimaryDeep}
\color{white}\bfseries متفکر & \color{white}\bfseries ایده‌ی اصلی
  & \color{white}\bfseries تأثیر بر انقلاب \\
\midrule
\rowcolor{BgWarm} \thinker{ولتر} & تساهل، نقد کلیسا & ضد روحانیت \\
\rowcolor{BgCool} \thinker{مونتسکیو} & تفکیک قوا & الهام‌بخش قانون اساسی \\
\rowcolor{BgWarm} \thinker{روسو} & قرارداد اجتماعی، اراده‌ی عمومی
  & دموکراسی رادیکال \\
\rowcolor{BgCool} \thinker{دیدرو} & دایرةالمعارف & نقد سنت و خرافه \\
\bottomrule
\end{tabular}
\end{adjustbox}
\end{center}

\begin{tcolorbox}[defbox, title={\rl{ایده‌های انقلابی}}]
\begin{itemize}[nosep, rightmargin=1em]
  \item \concept{حاکمیت ملی} به جای حق الهی شاهان
  \item \concept{برابری} در برابر قانون
  \item \concept{آزادی} فردی و اقتصادی
  \item \concept{عقل} به جای سنت و مذهب
\end{itemize}
\end{tcolorbox}

% ─── ۱.۴ بحران اجتماعی-اقتصادی ────────────────────────────
\needspace{6\baselineskip}%
\subsection{بحران اجتماعی–اقتصادی (۱۷۸۸–۸۹)}
\label{sec:social-crisis}

\begin{tcolorbox}[enemybox, title={\rl{قحطی و گرانی}}]
\begin{center}
\begin{adjustbox}{max width=\linewidth}%
\begin{tabular}{rr}
\toprule
\rowcolor{AccentRed!15}
\color{AccentRed}\bfseries شاخص & \color{AccentRed}\bfseries ۱۷۸۸–۸۹ \\
\midrule
قیمت نان & $+$۸۸٪ \\
درصد درآمد صرف نان (کارگر) & ۸۰–۹۰٪ \\
بیکاری در پاریس & ${\sim}$۵۰٪ \\
سرمای زمستان ۱۷۸۸ & شدیدترین در ۸۰ سال \\
\bottomrule
\end{tabular}
\end{adjustbox}
\end{center}
\end{tcolorbox}

% ═══════════════════════════════════════════════════════════════════
%   بخش دوم: فازهای انقلاب
% ═══════════════════════════════════════════════════════════════════
\newpage
\needspace{6\baselineskip}%
\section{فازهای انقلاب}
\label{sec:phases}

\subsection{فاز اول: انقلاب لیبرال–مشروطه (۱۷۸۹–۱۷۹۱)}
\label{sec:phase1}

\subsubsection{رویدادهای کلیدی}

\begin{center}
\begin{adjustbox}{max width=\linewidth}%
\begin{tabular}{rrr}
\toprule
\rowcolor{PhaseLiberal}
\color{white}\bfseries تاریخ & \color{white}\bfseries رویداد
  & \color{white}\bfseries اهمیت \\
\midrule
\rowcolor{BgCool} ۵ مه ۱۷۸۹ & گشایش مجلس طبقات & آغاز رسمی \\
\rowcolor{BgWarm} ۱۷ ژوئن ۱۷۸۹ & اعلام مجلس ملی & ادعای حاکمیت \\
\rowcolor{BgCool} ۲۰ ژوئن ۱۷۸۹ & سوگند زمین تنیس & تعهد به قانون اساسی \\
\rowcolor{AccentGold!20} \bfseries ۱۴ ژوئیه ۱۷۸۹ & \bfseries سقوط باستیل
  & \bfseries نماد انقلاب \\
\rowcolor{BgCool} ۴ اوت ۱۷۸۹ & لغو امتیازات فئودالی & پایان نظام کهن \\
\rowcolor{BgWarm} ۲۶ اوت ۱۷۸۹ & اعلامیه‌ی حقوق بشر & بنیاد ایدئولوژیک \\
\rowcolor{BgCool} ۵–۶ اکتبر ۱۷۸۹ & راهپیمایی زنان به ورسای
  & انتقال شاه به پاریس \\
\rowcolor{BgWarm} ۲۰–۲۱ ژوئن ۱۷۹۱ & فرار شاه (وارن)
  & بحران مشروعیت سلطنت \\
\rowcolor{BgCool} ۳ سپتامبر ۱۷۹۱ & قانون اساسی ۱۷۹۱ & سلطنت مشروطه \\
\bottomrule
\end{tabular}
\end{adjustbox}
\end{center}

\subsubsection{اعلامیه‌ی حقوق بشر و شهروند (۲۶ اوت ۱۷۸۹)}

\begin{tcolorbox}[mybox, title={\rl{اصول کلیدی اعلامیه‌ی حقوق بشر}}]
\begin{center}
\begin{adjustbox}{max width=\linewidth}%
\begin{tabular}{rr}
\toprule
\rowcolor{AccentGold!30}
\bfseries ماده & \bfseries محتوا \\
\midrule
\rowcolor{BgWarm} ۱ & انسان‌ها آزاد و برابر به دنیا می‌آیند \\
\rowcolor{BgCool} ۲ & حقوق طبیعی: آزادی، مالکیت، امنیت، مقاومت در برابر ستم \\
\rowcolor{BgWarm} ۳ & حاکمیت متعلق به ملت است \\
\rowcolor{BgCool} ۶ & قانون بیان اراده‌ی عمومی است \\
\rowcolor{BgWarm} ۱۰ & آزادی عقیده (حتی مذهبی) \\
\rowcolor{BgCool} ۱۱ & آزادی بیان \\
\rowcolor{BgWarm} ۱۷ & مالکیت حق مقدس است \\
\bottomrule
\end{tabular}
\end{adjustbox}
\end{center}
\end{tcolorbox}

\subsubsection{ساختار قانون اساسی ۱۷۹۱}

\begin{center}
\begin{tikzpicture}[every node/.style={font=\small, align=center},
  block/.style={rectangle, rounded corners=5pt, draw=PhaseLiberal,
    fill=PhaseLiberal!10, minimum width=4cm, minimum height=0.8cm},
  arr/.style={-{Stealth[length=5pt]}, thick, PhaseLiberal},
]
  \node[block] (king) {\rl{شاه (قوه‌ی مجریه)}\\[1pt]\rl{حق وتوی تعلیقی}};
  \node[block, below=0.8cm of king] (assembly)
    {\rl{مجلس قانون‌گذاری (قوه‌ی مقننه)}\\[1pt]\rl{۷۴۵ نماینده}};
  \node[block, below=0.8cm of assembly] (citizens)
    {\rl{شهروندان «فعال»}\\[1pt]\rl{~۴٫۳ میلیون مرد بالای ۲۵ سال}};
  \draw[arr] (king) -- (assembly);
  \draw[arr] (assembly) -- (citizens);
\end{tikzpicture}
\end{center}

\subsubsection{نیروهای سیاسی}

\begin{center}
\begin{adjustbox}{max width=\linewidth}%
\begin{tabular}{rrr}
\toprule
\rowcolor{PrimaryDeep}
\color{white}\bfseries گروه & \color{white}\bfseries موضع
  & \color{white}\bfseries رهبران \\
\midrule
\rowcolor{BgCool} فویان‌ها & سلطنت مشروطه، میانه‌رو & لافایت، بایی \\
\rowcolor{BgWarm} ژیروندن‌ها & جمهوری، جنگ & بریسو، رولان \\
\rowcolor{PhaseRadical!10} ژاکوبن‌ها & جمهوری رادیکال & روبسپیر، مارا \\
\rowcolor{BgCool} کوردلیه‌ها & دموکراسی مستقیم & دانتون، مارا \\
\bottomrule
\end{tabular}
\end{adjustbox}
\end{center}

% ─── فاز دوم ─────────────────────────────────────────────────
\newpage
\needspace{6\baselineskip}%
\subsection{فاز دوم: انقلاب رادیکال و ترور (۱۷۹۲–۱۷۹۴)}
\label{sec:phase2}

\subsubsection{از سلطنت به جمهوری}

\begin{center}
\begin{adjustbox}{max width=\linewidth}%
\begin{tabular}{rrr}
\toprule
\rowcolor{PhaseRadical}
\color{white}\bfseries تاریخ & \color{white}\bfseries رویداد
  & \color{white}\bfseries اهمیت \\
\midrule
\rowcolor{BgWarm} ۲۰ آوریل ۱۷۹۲ & اعلان جنگ به اتریش & آغاز جنگ‌های انقلابی \\
\rowcolor{AccentRed!10} \bfseries ۱۰ اوت ۱۷۹۲ & \bfseries حمله به تویلری
  & \bfseries سقوط سلطنت \\
\rowcolor{BgWarm} ۲–۶ سپتامبر ۱۷۹۲ & کشتارهای سپتامبر
  & ترور خودجوش (~۱٬۴۰۰ کشته) \\
\rowcolor{BgCool} ۲۱ سپتامبر ۱۷۹۲ & اعلام جمهوری & جمهوری اول \\
\rowcolor{AccentRed!15} \bfseries ۲۱ ژانویه ۱۷۹۳ & \bfseries اعدام لویی شانزدهم
  & \bfseries نقطه‌ی بازگشت‌ناپذیر \\
\bottomrule
\end{tabular}
\end{adjustbox}
\end{center}

\subsubsection{ساختار کنوانسیون ملی}

\begin{center}
\begin{tikzpicture}[every node/.style={font=\small, align=center},
  block/.style={rectangle, rounded corners=5pt, draw=#1, fill=#1!10,
    minimum width=2.5cm, minimum height=1cm, line width=1pt},
  arr/.style={-{Stealth[length=4pt]}, thick, NeutralDark},
]
  \node[block=PrimaryDeep] (conv)
    {\rl{کنوانسیون ملی}\\[1pt]\rl{۷۴۹ نماینده}};
  \node[block=PhaseRadical, below right=1.5cm and -0.5cm of conv] (mont)
    {\rl{کوه (مونتانیار)}\\[1pt]\rl{~۳۰۰ رادیکال}};
  \node[block=NeutralMid, below=1.5cm of conv] (plain)
    {\rl{دشت (مرداب)}\\[1pt]\rl{~۳۵۰ میانه}};
  \node[block=PhaseLiberal, below left=1.5cm and -0.5cm of conv] (gir)
    {\rl{ژیروند}\\[1pt]\rl{~۱۶۰ میانه‌رو}};
  \draw[arr] (conv) -- (mont);
  \draw[arr] (conv) -- (plain);
  \draw[arr] (conv) -- (gir);
\end{tikzpicture}
\end{center}

\subsubsection{دوران ترور}

\begin{tcolorbox}[enemybox, title={\rl{دوران ترور: ژوئن ۱۷۹۳ – ژوئیه ۱۷۹۴}}]

\begin{center}
\begin{adjustbox}{max width=\linewidth}%
\begin{tabular}{rr}
\toprule
\rowcolor{AccentRed!15}
\color{AccentRed}\bfseries تاریخ & \color{AccentRed}\bfseries رویداد \\
\midrule
\rowcolor{BgWarm} ۲ ژوئن ۱۷۹۳ & سقوط ژیروندن‌ها \\
\rowcolor{BgCool} ۱۳ ژوئیه ۱۷۹۳ & ترور مارا (توسط شارلوت کوردی) \\
\rowcolor{BgWarm} ۵ سپتامبر ۱۷۹۳ & «ترور در دستور کار» \\
\rowcolor{BgCool} ۱۶ اکتبر ۱۷۹۳ & اعدام ماری آنتوانت \\
\rowcolor{BgWarm} ۳۱ اکتبر ۱۷۹۳ & اعدام ژیروندن‌ها \\
\rowcolor{BgCool} ۵ آوریل ۱۷۹۴ & اعدام دانتون و میانه‌روها \\
\rowcolor{BgWarm} ۱۰ ژوئن ۱۷۹۴ & قانون ۲۲ پریریال (تسریع ترور) \\
\rowcolor{AccentRed!20} \bfseries ۲۷ ژوئیه ۱۷۹۴
  & \bfseries ۹ ترمیدور: سقوط روبسپیر \\
\bottomrule
\end{tabular}
\end{adjustbox}
\end{center}

\vspace{0.5\baselineskip}

\textbf{آمار ترور:}
\begin{center}
\begin{adjustbox}{max width=0.7\linewidth}%
\begin{tabular}{rr}
\toprule
\bfseries شاخص & \bfseries تعداد \\
\midrule
اعدام با گیوتین & ${\sim}$۱۷٬۰۰۰ \\
مرگ در زندان & ${\sim}$۱۰٬۰۰۰ \\
کشته در سرکوب واندی & ${\sim}$۱۵۰٬۰۰۰–۲۵۰٬۰۰۰ \\
\midrule
\bfseries مجموع & \bfseries ${\sim}$۲۰۰٬۰۰۰–۳۰۰٬۰۰۰ \\
\bottomrule
\end{tabular}
\end{adjustbox}
\end{center}
\end{tcolorbox}

\begin{tcolorbox}[quotebox]
«ترور چیزی نیست جز عدالت سریع، شدید و انعطاف‌ناپذیر؛
پس جلوه‌ای از فضیلت است.»

\hfill — روبسپیر، ۵ فوریه ۱۷۹۴
\end{tcolorbox}

\subsubsection{نهادهای ترور}

\begin{center}
\begin{adjustbox}{max width=\linewidth}%
\begin{tabular}{rr}
\toprule
\rowcolor{AccentRed!15}
\bfseries نهاد & \bfseries نقش \\
\midrule
\rowcolor{BgWarm} کمیته‌ی نجات عمومی & قوه‌ی مجریه، هدایت جنگ \\
\rowcolor{BgCool} کمیته‌ی امنیت عمومی & پلیس سیاسی \\
\rowcolor{BgWarm} دادگاه انقلابی & محاکمه‌ی سریع \\
\rowcolor{BgCool} نمایندگان مأمور & اجرای ترور در ولایات \\
\bottomrule
\end{tabular}
\end{adjustbox}
\end{center}

% ─── فاز سوم ─────────────────────────────────────────────────
\newpage
\needspace{6\baselineskip}%
\subsection{فاز سوم: ترمیدور و دیرکتوار (۱۷۹۴–۱۷۹۹)}
\label{sec:phase3}

\subsubsection{واکنش ترمیدوری}

پس از سقوط روبسپیر، موجی از واکنش علیه رادیکالیسم آغاز شد.
«ترور سفید» جای ترور سرخ را گرفت و نهایتاً قانون اساسی جدیدی
نوشته شد که سعی داشت میان آزادی و نظم تعادل برقرار کند.

\begin{center}
\begin{adjustbox}{max width=\linewidth}%
\begin{tabular}{rr}
\toprule
\rowcolor{PhaseThermidor}
\color{white}\bfseries تاریخ & \color{white}\bfseries رویداد \\
\midrule
\rowcolor{BgWarm} ۲۷ ژوئیه ۱۷۹۴ & کودتای ۹ ترمیدور \\
\rowcolor{BgCool} ۲۸ ژوئیه ۱۷۹۴ & اعدام روبسپیر و یارانش \\
\rowcolor{BgWarm} ۱۷۹۴–۹۵ & «ترور سفید» علیه ژاکوبن‌ها \\
\rowcolor{BgCool} ۲۲ اوت ۱۷۹۵ & قانون اساسی سال سوم \\
\rowcolor{BgWarm} ۲۶ اکتبر ۱۷۹۵ & آغاز دیرکتوار \\
\bottomrule
\end{tabular}
\end{adjustbox}
\end{center}

\subsubsection{ساختار دیرکتوار}

\begin{center}
\begin{tikzpicture}[every node/.style={font=\small, align=center},
  block/.style={rectangle, rounded corners=5pt, draw=PhaseThermidor,
    fill=PhaseThermidor!10, minimum width=4.5cm, minimum height=0.8cm},
  arr/.style={-{Stealth[length=5pt]}, thick, PhaseThermidor},
]
  \node[block] (dir) {\rl{۵ دیرکتور (قوه‌ی مجریه)}};
  \node[block, below=0.7cm of dir] (coun)
    {\rl{شورای ۵۰۰ + شورای ریش‌سفیدان}};
  \node[block, below=0.7cm of coun] (elec)
    {\rl{انتخابات محدود (مالکین)}};
  \draw[arr] (dir) -- (coun);
  \draw[arr] (coun) -- (elec);
\end{tikzpicture}
\end{center}

\begin{tcolorbox}[defbox, title={\rl{مشکلات دیرکتوار}}]
\begin{itemize}[nosep, rightmargin=1em]
  \item بی‌ثباتی سیاسی (کودتاهای مکرر)
  \item فساد اقتصادی گسترده
  \item وابستگی فزاینده به ارتش
  \item تهدید همزمان سلطنت‌طلبان و ژاکوبن‌ها
\end{itemize}
\end{tcolorbox}

\subsubsection{کودتای ۱۸ برومر (۹ نوامبر ۱۷۹۹)}

\begin{center}
\begin{tikzpicture}[every node/.style={font=\small, align=center},
  block/.style={rectangle, rounded corners=5pt, draw=PhaseNapoleon,
    fill=PhaseNapoleon!10, minimum width=4cm, minimum height=0.8cm},
  arr/.style={-{Stealth[length=5pt]}, thick, PhaseNapoleon},
]
  \node[block] (coup) {\rl{ناپلئون + سیئس + دوکو}};
  \node[block, below=0.6cm of coup] (diss) {\rl{انحلال دیرکتوار}};
  \node[block, below=0.6cm of diss] (cons) {\rl{کنسولا (سه کنسول)}};
  \node[block, below=0.6cm of cons] (nap) {\rl{ناپلئون: کنسول اول}};
  \draw[arr] (coup) -- (diss);
  \draw[arr] (diss) -- (cons);
  \draw[arr] (cons) -- (nap);
\end{tikzpicture}
\end{center}

% ═══════════════════════════════════════════════════════════════════
%   بخش سوم: دوران ناپلئون
% ═══════════════════════════════════════════════════════════════════
\newpage
\needspace{6\baselineskip}%
\section{دوران ناپلئون (۱۷۹۹–۱۸۱۵)}
\label{sec:napoleon}

\subsection{کنسولا (۱۷۹۹–۱۸۰۴)}
\label{sec:consulate}

\subsubsection{اصلاحات ناپلئونی}

ناپلئون با ترکیبی از اقتدارگرایی و مدرن‌سازی، بسیاری از دستاوردهای
انقلاب را نهادینه کرد و ساختارهایی ایجاد کرد که تا امروز پابرجاست.

\begin{center}
\begin{adjustbox}{max width=\linewidth}%
\begin{tabular}{rrrr}
\toprule
\rowcolor{PhaseNapoleon}
\color{white}\bfseries حوزه & \color{white}\bfseries اصلاح
  & \color{white}\bfseries سال & \color{white}\bfseries اهمیت \\
\midrule
\rowcolor{BgWarm} حقوقی & کد ناپلئون & ۱۸۰۴ & بنیاد حقوق مدرن \\
\rowcolor{BgCool} اداری & پرفه‌ها (استانداران) & ۱۸۰۰ & تمرکز اداری \\
\rowcolor{BgWarm} مالی & بانک فرانسه & ۱۸۰۰ & ثبات مالی \\
\rowcolor{BgCool} آموزشی & لیسه‌ها & ۱۸۰۲ & آموزش متمرکز \\
\rowcolor{BgWarm} مذهبی & کنکوردا با پاپ & ۱۸۰۱ & آشتی با کلیسا \\
\bottomrule
\end{tabular}
\end{adjustbox}
\end{center}

\subsubsection{کد ناپلئون (۱۸۰۴)}

\begin{tcolorbox}[wavebox, title={\rl{کد ناپلئون: اصول و محدودیت‌ها}}]
\textbf{اصول کلیدی:}
\begin{itemize}[nosep, rightmargin=1em]
  \item برابری در برابر قانون
  \item حق مالکیت مقدس
  \item سکولاریسم
  \item آزادی قرارداد
\end{itemize}

\vspace{0.3\baselineskip}
\textbf{محدودیت‌ها:}
\begin{itemize}[nosep, rightmargin=1em]
  \item زنان تحت قیمومیت شوهر
  \item حقوق کارگران محدود
  \item برده‌داری در مستعمرات بازگردانده شد
\end{itemize}
\end{tcolorbox}

% ─── ۳.۲ امپراتوری ────────────────────────────────────────────
\needspace{6\baselineskip}%
\subsection{امپراتوری اول (۱۸۰۴–۱۸۱۵)}
\label{sec:empire1

ادامه‌ی فایل لاتکس: